\documentclass[a4paper,10pt,notitlepage]{article}
\usepackage{ctex,geometry,graphicx,tikz,setspace,paralist,fancyhdr,caption}
\geometry{
	left=1.5cm,
	right=1.5cm,
	top=1.4cm,
	bottom=2cm,}
\newcommand{\rec}{
	\begin{tikzpicture}[remember picture,overlay]
		% 绘制边框
		\draw[line width=1.2pt] ([xshift=0.5cm,yshift=0.5cm] current page.south west) rectangle ([xshift=-0.5cm,yshift=-0.8cm] current page.north east);
	\end{tikzpicture}
	
}
\pagestyle{fancy}
\fancyhf{}
\fancyhead[C]{\rec}  % 在页眉中绘制图形
\fancyhead[L]{附录}
\fancyhead[R]{实验八 \quad 波形发生电路}
\fancyfoot[C]{\thepage}
\begin{document}
	\large
	\onehalfspacing
	\begin{figure}[h]
		\raggedright
		\includegraphics{1.png}
	\end{figure}
	\centering
	{\Huge\textbf{模电实验报告}\par}
	\vspace{0.2cm}
	{\huge{实验内容：波形发生电路}\par}
	\raggedright
	\vspace{0.3cm}
	\begin{centering}
		{\large 院系：电子与信息工程学院\hfill 学号：22309080\hfill 审批：\hspace{2cm} \par
			专业：通信工程\hfill 实验人：梁倍铭\hfill 日期：2023年11月30日 \par}
	\end{centering}
	\vspace{0.3cm}
	\section*{一、方波发生电路}
		仿真波形如图
		\begin{figure}[h]
			\centering
			\begin{minipage}{0.3\textwidth}
				\centering
				\includegraphics[width=\textwidth]{2.png}
				\caption*{图1 波形输出}
			\end{minipage}
			\qquad
			\begin{minipage}{0.3\textwidth}
				\centering
				\begin{tabular}{|c|c|c|}
					\hline
					 & 频率(Hz) & 幅值(V) \\
					\hline
					仿真 & 388.8 & 6.47 \\
					\hline 
					实验 & \qquad & \qquad  \\
					\hline
				\end{tabular}
				\caption*{表1 10k时的频率与幅值}
			\end{minipage}
			\qquad
			\begin{minipage}{0.3\textwidth}
				\centering
				\begin{tabular}{|c|c|c|}
					\hline
					& 频率(Hz) & 幅值(V) \\
					\hline
					仿真 & 40.66 & 8.55 \\
					\hline 
					实验 & \qquad & \qquad  \\
					\hline
				\end{tabular}
				\caption*{表2 110k时的频率与幅值}
			\end{minipage}
		\end{figure}
		想要获得更低频率增大R即可
	\section*{二、占空比可调的矩形波发生电路}
	\begin{figure}[h]
		\centering
		\begin{minipage}{0.3\textwidth}
			\centering
			\includegraphics[width=\textwidth]{3.png}
			\caption*{图2 波形输出}
		\end{minipage}
		\qquad
		\begin{minipage}{0.3\textwidth}
			\centering
			\begin{tabular}{|c|c|c|c|}
				\hline
				& 频率(Hz) & 幅值(V) & 占空比(\%) \\
				\hline
				仿真 & 77.52 & 8.348 & 49.3 \\
				\hline 
				实验 & \qquad & \qquad & \qquad \\
				\hline
			\end{tabular}
			\caption*{表3 观测数据}
		\end{minipage}
	\end{figure}
	想要波形占空比更大，调节RP1使得正向充电的时间更长即可
	\newpage
	\section*{三、三角波发生电路}
	\begin{figure}[h]
		\centering
		\includegraphics[width=0.3\textwidth]{4.png}
		\caption*{图3 波形输出}
	\end{figure}
	可以通过改变RP的阻值来改变波形频率
	\section*{四、锯齿波发生电路}
	\begin{figure}[h]
		\centering
		\includegraphics[width=0.3\textwidth]{5.png}
		\caption*{图4 波形输出}
	\end{figure}
可以通过改变滑动变阻器来改变频率变化范围，仿真结果为：24.57Hz--40.05Hz,实验结果为：
\end{document}